\subsection*{Поиск порядковой статистики}

\Alg

Для поиска порядковой статистики будем хранить вспомогательную величину Size(): количество вершин в поддереве нашей вершины (в поддерево включается и сама вершина).

\begin{enumerate}
    \item left.Size() = k, то завершаем поиск
    \item left.Size() > k, то рекурсивно ищем в левом поддереве k
    \item left.Size() < k, то рекурсивно ищем в правом поддереве (k - left.Size() - 1)
\end{enumerate}

\begin{lstlisting}
optional<T> FindKthElement(Node* node, size_t serial) {
  if (node == nullptr || serial > Size(node)) {
    return {};
  }
  if (Size(node->left) == serial) {
    return node->value;
  }
  if (Size(node->left) > serial) {
    return FindKthElement(node->left, serial);
  }
  return FindKthElement(node->right, serial - Size(node->left) - 1);
}
\end{lstlisting}

\subsubsection*{LowerBound}

\Def LowerBound - поиск первого элемента, который не больше чем данное значение.

\Alg

\begin{enumerate}
    \item Если заданное значение равно node.value, то завершаем поиск и возвращаем node.value.
    \item Если значение в вершине меньше чем заданное, то рекурсивно ищем в правом поддереве. Если мы нашли значение, то возвращаем его, иначе возвращаем значение текущей вершины.
    \item Если значение в вершине больше чем заданное, рекурсивно ищем в левом поддереве
\end{enumerate}

\begin{lstlisting}
optional<T> LowerBound(Node* node, const T& value) {
  if (node == nullptr) {
    return {};
  }
  if (node->value == value) {
    return node->value;
  }
  if (node->value < value) {
    optional<T> next = LowerBound(node->right, value);
    if (next.has_value()) {
      return next.value();
    }
    return node->value;
  }
  return LowerBound(node->left, value);
}
\end{lstlisting}

\Statement Операции FindKthElement, LowerBound работают за $O(h)$, где h -- высота дерева

\pagebreak